\documentclass[11pt,a4paper]{article}

\usepackage[utf8]{inputenc}
\usepackage[T1]{fontenc}
\usepackage[french]{babel}
\usepackage{lmodern}
\usepackage[a4paper,margin=2.4cm]{geometry}
\usepackage{graphicx}
\usepackage{booktabs}
\usepackage{tabularx}
\usepackage{array}
\usepackage{enumitem}
\usepackage{caption}
\usepackage{float}
\usepackage{tikz}
\usetikzlibrary{arrows.meta,positioning,fit,calc}
\usepackage[hidelinks]{hyperref}

\graphicspath{{images/}{./}}
\setlist{itemsep=2pt,topsep=4pt}
\setlength{\parskip}{5pt}
\setlength{\parindent}{0pt}
\setlength{\emergencystretch}{2em}
\renewcommand{\arraystretch}{1.25}
\newcolumntype{L}{>{\raggedright\arraybackslash}X}

\newcommand{\code}[1]{\texttt{#1}}

% Insere la capture si le fichier existe, sinon un cadre de remplacement
\newcommand{\capture}[3]{%
  \begin{figure}[H]
    \centering
    \IfFileExists{images/#1}{%
      \includegraphics[width=\linewidth,height=0.4\textheight,keepaspectratio]{#1}%
    }{%
      \IfFileExists{#1}{%
        \includegraphics[width=\linewidth,height=0.4\textheight,keepaspectratio]{#1}%
      }{%
        \fbox{\parbox[c][5.5cm][c]{0.92\linewidth}{\centering
          \textit{Capture à insérer :} \code{#1}\\[4pt]
          \small #2}}%
      }%
    }
    \caption{#2}
    \label{#3}
  \end{figure}
}

\begin{document}

% ------------------------------------------------------------------
\begin{titlepage}
  \centering
  \vspace*{2cm}
  {\large Université Côte d'Azur\par}
  {\large Master 2 Intelligence Artificielle\par}
  \vspace{2.5cm}
  {\large Projet : Multi-Agent Systems\par}
  \vspace{0.8cm}
  {\LARGE\bfseries Conception d'un système multi-agent\\[4pt]
   pour la réponse automatique aux emails\\[4pt]
   à l'aide de LLM locaux\par}
  \vspace{2.5cm}
  {\large Mounia AREZZOUG --- 22509499\par}
  \vspace{0.3cm}
  {\large Yunus Emre DOGAN --- 22508757\par}
  \vspace{0.3cm}
  {\large Mathis LEVESQUE DU ROSTU --- 22510354\par}
  \vfill
  {Année universitaire 2026--2027\par}
\end{titlepage}

\tableofcontents
\newpage

% ------------------------------------------------------------------
\section{Introduction et problématique}

Un enseignant-chercheur reçoit des emails très différents : questions d'étudiants,
messages de collègues, demandes de l'administration, échanges autour d'un projet de
recherche. Une bonne réponse dépend de l'expéditeur, du sujet, de l'historique de la
conversation et d'informations propres à chaque contexte.

Les grands modèles de langage savent rédiger ce type de réponse, mais les services en
ligne posent un problème de confidentialité : ces emails contiennent des données
personnelles ou non publiées. Le sujet impose donc d'utiliser un LLM exécuté
\emph{localement}. Ces modèles sont plus petits et moins fiables que les modèles distants.

La question posée par le sujet est la suivante : \emph{une organisation d'agents
spécialisés permet-elle d'utiliser efficacement des LLM locaux de taille raisonnable pour
produire des réponses adaptées à différents contextes ?}

Notre réponse est un programme Java qui ne fait pas « tout passer » dans un seul modèle :
un agent de routage choisit l'agent spécialisé le plus adapté. Cet agent rédige un
brouillon avec son propre contexte et un agent critique le relit. L'utilisateur garde
toujours la décision finale. Aucun email n'est envoyé automatiquement.

% ------------------------------------------------------------------
\section{Notions de systèmes multi-agents mobilisées}

\paragraph{Qu'est-ce qu'un agent ?}
Un agent est une entité \emph{autonome} qui perçoit son environnement, possède ses propres
connaissances et agit pour atteindre un objectif. Dans notre projet, un agent spécialisé
reçoit un email (sa perception). Il dispose d'un rôle, d'un \emph{system prompt}, d'un
contexte qui lui est propre et d'une mémoire des conversations, et il agit en proposant
une réponse. Deux agents différents ne savent pas les mêmes choses : l'agent
administratif ne voit pas les fichiers du projet de recherche.

\paragraph{Le modèle BDI.}
Nos agents suivent le modèle \emph{Belief--Desire--Intention} :
\begin{itemize}
  \item \textbf{Croyances} : les faits du contexte (\code{date | fait}) et les échanges validés des fils déjà traités ;
  \item \textbf{Désir} : obtenir une réponse validée par l'utilisateur en un nombre limité d'essais (\code{maxRetries}) ;
  \item \textbf{Intention} : l'engagement concret à répondre à \emph{cet} email, avec ses brouillons, les refus et leurs raisons ;
  \item \textbf{Plan} : rédiger, soumettre, puis mémoriser la réponse, recommencer avec les remarques ou abandonner.
\end{itemize}

\paragraph{Pourquoi plusieurs agents plutôt qu'un seul programme ?}
Avec un seul agent, il faudrait envoyer au modèle toutes les règles et toutes les
informations de tous les contextes. C'est long, coûteux et source d'erreurs pour un petit
modèle. En spécialisant les agents, chaque appel au LLM reçoit un prompt court et ciblé.
Les contextes sont isolés les uns des autres, ce qui limite aussi les fuites
d'informations. Enfin, l'organisation est extensible : ajouter un domaine revient à
ajouter un agent.

\paragraph{Organisation et coordination.}
Notre organisation est hiérarchique et simple. Un agent coordinateur (le
\code{RoutingAgent}) reçoit les emails, choisit l'agent compétent et sert
d'intermédiaire entre cet agent, l'agent critique et l'utilisateur. Les agents
spécialisés ne communiquent jamais directement entre eux : toute interaction passe par le
coordinateur, qui joue le rôle de \emph{médiateur} (y compris quand plusieurs agents
collaborent sur un email, section~\ref{sec:collab}). La communication se fait par
échange d'objets-messages explicites (email, intention, proposition, retour de
l'utilisateur, critique).

% ------------------------------------------------------------------
\section{Architecture du système}

\subsection{Vue d'ensemble}

La figure~\ref{fig:archi} présente l'organisation des agents. Le code est écrit en Java
17, sans framework multi-agent externe. La seule bibliothèque utilisée est Jackson, pour
le JSON.

\begin{figure}[H]
\centering
\begin{tikzpicture}[
  font=\small,
  box/.style={draw, rounded corners=2pt, minimum height=8mm, align=center, inner sep=3pt},
  agent/.style={box, minimum width=35mm, minimum height=6.5mm},
  arr/.style={-{Stealth[length=2.2mm]}, thick},
  darr/.style={-{Stealth[length=2.2mm]}, thick, dashed}
]
\node[box, minimum width=24mm] (user) at (0,0.2) {Utilisateur\\(interface Swing)};
\node[box, minimum width=25mm, line width=1pt] (router) at (4.0,0.2) {RoutingAgent\\(coordinateur)};
\node[agent] (s) at (9.1,2.0) {StudentAgent};
\node[agent] (t) at (9.1,1.1) {TeacherAgent};
\node[agent] (a) at (9.1,0.2) {AdministrationAgent};
\node[agent] (r) at (9.1,-0.7) {ResearchProjectAgent};
\node[agent] (g) at (9.1,-1.6) {GenericAgent};
\node[draw, dashed, rounded corners, fit=(s)(g), inner sep=4pt, label=above:{Agents spécialisés (BDI)}] (spec) {};
\node[box, minimum width=23mm] (llm) at (13.4,0.2) {LLM local\\(serveur\\compatible\\OpenAI)};
\node[box, minimum width=33mm] (signals) at (4.0,-2.0) {Signaux de routage\\{\footnotesize Thread, Contact, Keyword, Domain}};
\node[box, minimum width=33mm] (arbiter) at (4.0,-3.6) {Arbitre LLM\\(LlmClassifierStep)};
\node[agent, minimum height=8mm] (critic) at (9.1,-3.6) {CriticAgent};

\draw[arr] (user) -- node[above]{email} (router);
\draw[arr] (router.north) to[bend right=30] node[above]{proposition} (user.north);
\draw[arr] (router) -- node[above, pos=0.4]{intention} (spec.west |- router);
\draw[arr] (router) -- (signals);
\draw[darr] (signals) -- node[right]{\footnotesize si doute} (arbiter);
\draw[arr] (6.55,0.2) -- node[right, pos=0.72]{\footnotesize brouillon} (6.55,-3.6) -- (critic.west);
\draw[arr] (spec.east |- llm) -- (llm);
\draw[arr] (critic.east) -| (llm.south);
\draw[arr] (arbiter.south) -- ++(0,-0.45) -| ($(llm.south)+(0.6,0)$);
\end{tikzpicture}
\caption{Organisation des agents. Toutes les interactions passent par le RoutingAgent, qui peut solliciter plusieurs agents spécialisés sur un email ambigu (section~\ref{sec:collab}).}
\label{fig:archi}
\end{figure}

\subsection{Les agents et leurs responsabilités}

\begin{table}[H]
\centering
\small
\begin{tabularx}{\linewidth}{l L l}
\toprule
\textbf{Agent} & \textbf{Responsabilité} & \textbf{Contexte propre} \\
\midrule
RoutingAgent & Choisit l'agent compétent (et, si l'email est ambigu, des agents contributeurs), transmet l'email, gère la boucle de validation. & contacts, mots-clés \\
StudentAgent & Répond aux étudiants : cours, TP, rendus, examens. & 10 faits \\
TeacherAgent & Répond aux collègues : réunions, jurys, enseignement. & 10 faits \\
AdministrationAgent & Répond à l'administration : documents, signatures, échéances, rendez-vous (ton formel). & 10 faits \\
ResearchProjectAgent & Gère le projet de recherche \emph{SuperMASProject} : membres, livrables, deadlines. & 2 fichiers, 14 faits \\
GenericAgent & Agent par défaut quand aucun autre ne convient. & 2 faits \\
CriticAgent & Évalue chaque brouillon (note, remarques). N'est jamais choisi par le routeur. & celui de l'agent relu \\
\bottomrule
\end{tabularx}
\caption{Agents présents dans le système.}
\end{table}

Chaque agent spécialisé est défini par un dossier \code{agents/<nom>/} contenant un fichier
\code{agent.json} (nom, rôle, description, mots-clés, \code{maxRetries}), un
\emph{system prompt} et un ou plusieurs fichiers de contexte. La classe \code{AgentFactory}
crée automatiquement un \code{BDIAgent} pour chaque dossier. Ajouter un agent ne demande
donc aucune modification du code Java. L'\code{AdministrationAgent} et le
\code{ResearchProjectAgent} ont d'ailleurs été ajoutés de cette façon.

Le \emph{system prompt} définit le rôle, le ton, ce que l'agent doit faire et ce qu'il ne
doit jamais faire. Par exemple, le \code{ResearchProjectAgent} doit toujours rappeler la
prochaine échéance, ne jamais s'engager sur une nouvelle tâche sans validation et ne
jamais partager de résultats non publiés avec une personne extérieure au projet.

\subsection{Principales classes et design patterns}

\begin{table}[H]
\centering
\small
\begin{tabularx}{\linewidth}{l l L}
\toprule
\textbf{Pattern} & \textbf{Où} & \textbf{Pourquoi} \\
\midrule
Factory & \code{AgentFactory} & Créer un agent à partir de sa configuration, sans code spécifique. \\
Strategy & \code{RoutingStep}, \code{Arbiter} & Chaque signal de routage est interchangeable ; en ajouter un revient à écrire une classe. \\
Adapter & \code{LmStudioClient} & Adapte l'API HTTP du serveur local à une seule méthode \code{ask(prompt, message)}. \\
Mediator & \code{RoutingAgent} & Les agents, le critique et l'utilisateur ne se connaissent pas entre eux. \\
Observer & \code{WorkflowListener} & Le routeur publie ses étapes ; l'interface les affiche sans que le routeur la connaisse. \\
Facade & \code{EmailAssistant} & Construit et cache l'assemblage du système pour l'interface et les tests. \\
Repository & \code{AgentRegistry} & Lit et modifie les fichiers des agents depuis l'interface. \\
Strategy & \code{ContextSelector} & Choisir comment sélectionner le contexte (tout, BM25, embeddings). \\
Factory & \code{ContextSelectorFactory} & Créer la stratégie demandée par la configuration. \\
Decorator & \code{CachingEmbedder}, \code{MeteredLLM} & Ajouter un cache ou des mesures sans modifier le client. \\
Observer & \code{DecisionLogger} & Journaliser sans que le routeur connaisse le journal. \\
\bottomrule
\end{tabularx}
\caption{Design patterns réellement utilisés.}
\end{table}

L'interface \code{Validator} sépare la validation humaine du reste du système : la version
terminale (\code{ConsoleValidator}) et l'interface graphique (\code{SwingValidator})
l'implémentent sans changer la boucle du routeur.

% ------------------------------------------------------------------
\section{Organisation, communication et coordination}

\subsection{Le routage}

Le \code{RoutingAgent} ne s'arrête pas au premier indice. Quatre signaux indépendants
votent chacun pour un agent avec une confiance entre 0 et 1, et les votes sont
additionnés par agent.

\begin{table}[H]
\centering
\small
\begin{tabularx}{\linewidth}{l l L}
\toprule
\textbf{Signal} & \textbf{Vote} & \textbf{Condition} \\
\midrule
Fil (\code{ThreadStep}) & 0,9 & Un agent a déjà traité ce fil de discussion. \\
Contact (\code{ContactStep}) & 1,0 & L'expéditeur est connu dans \code{contacts.json}. \\
Mots-clés (\code{KeywordStep}) & 0,5 à 0,75 & L'agent qui a le plus de mots-clés dans le sujet et le corps (aucun vote en cas d'égalité). \\
Domaine (\code{DomainStep}) & 0,6 & Expéditeur inconnu : domaine étudiant, domaine universitaire ou boîte administrative (\code{scolarite@}, \code{rh@}...). \\
\bottomrule
\end{tabularx}
\caption{Signaux de routage.}
\end{table}

Un agent est choisi directement si son score représente au moins 75\,\% du total
\emph{et} vaut au moins 0,7. Un contact ou un fil connu suffisent donc seuls, alors que
des mots-clés ou un domaine doivent se confirmer mutuellement.

\paragraph{Plusieurs agents conviennent.} Si les votes sont partagés, le routeur consulte
le LLM comme arbitre (\code{LlmClassifierStep}). Il lui donne seulement les agents en
conflit et leurs descriptions, et le modèle répond en JSON avec un rôle et une confiance.
Il peut aussi désigner des agents contributeurs pour une partie distincte de l'email
(section~\ref{sec:collab}).

\paragraph{Aucun agent ne convient.} Si les signaux sont faibles, l'arbitre choisit parmi
tous les agents. Son choix n'est retenu que si sa confiance atteint 0,6 ; sinon,
l'email est confié au \code{GenericAgent}. La raison de la décision (liste des votes,
conflit, arbitrage) est conservée et affichée.

\subsection{Contexte et mémoire des conversations}

Seuls le contexte de l'agent choisi et l'historique du fil courant sont utilisés.
Parmi eux, seuls les faits les plus pertinents pour l'email et les derniers messages du
fil sont envoyés au modèle (section~\ref{sec:contexte}). Cela évite de surcharger un
petit LLM avec des informations inutiles.

Quand l'utilisateur valide une réponse, l'agent mémorise deux croyances : l'email reçu et
la réponse réellement retenue, éventuellement corrigée par l'utilisateur. Elles sont
sauvegardées dans \code{learnt\_beliefs\_<agent>.txt} et rechargées au démarrage. Les
brouillons refusés ne sont jamais mémorisés. Cette mémoire sert deux fois : le
\code{ThreadStep} renvoie une relance au même agent, et cet agent reçoit les derniers
messages du fil pour rédiger.

\subsection{Agent critique et validation humaine}

Avant d'être montré à l'utilisateur, chaque brouillon est relu par le \code{CriticAgent}.
Celui-ci reçoit les règles et le contexte de l'agent auteur, l'email et le brouillon. Il
renvoie une note sur 10 et des remarques : réponse incomplète, fait inventé, règle non
respectée. Un brouillon n'est considéré comme approuvé qu'à partir de 7/10. Le critique
conseille sans bloquer : s'il échoue, le traitement continue. Il peut aussi être désactivé.

L'utilisateur garde toujours la décision :
\begin{itemize}
  \item \textbf{accepter} la réponse, ou la \textbf{modifier} puis la valider ;
  \item \textbf{demander un nouveau brouillon} avec des remarques (éventuellement celles du critique) ;
  \item \textbf{refuser} définitivement : rien n'est mémorisé.
\end{itemize}
Après \code{maxRetries} refus, l'intention est abandonnée et l'utilisateur peut écrire la
réponse lui-même.

% ------------------------------------------------------------------
\section{Déroulement d'un scénario complet}

Prenons l'email d'un étudiant connu (Luca Martin) qui demande quels chapitres tombent à
l'examen. La figure~\ref{fig:sequence} résume les échanges.

\begin{figure}[H]
\centering
\begin{tikzpicture}[
  font=\small,
  etape/.style={draw, rounded corners=2pt, minimum width=10.5cm, minimum height=7mm, align=left, inner sep=4pt},
  arr/.style={-{Stealth[length=2mm]}, thick}
]
\node[etape] (e1) {1. L'utilisateur saisit l'email et lance le traitement.};
\node[etape, below=0.3cm of e1] (e2) {2. Le RoutingAgent collecte les votes : contact connu (1,0) + mot-clé « exam » (0,5).};
\node[etape, below=0.3cm of e2] (e3) {3. Votes concordants : StudentAgent choisi sans appel à l'arbitre.};
\node[etape, below=0.3cm of e3] (e4) {4. StudentAgent adopte une intention et récupère contexte + mémoire du fil.};
\node[etape, below=0.3cm of e4] (e5) {5. Il envoie au LLM local son prompt, son contexte et l'email : brouillon 1.};
\node[etape, below=0.3cm of e5] (e6) {6. Le CriticAgent note le brouillon et propose des remarques.};
\node[etape, below=0.3cm of e6] (e7) {7. Le RoutingAgent présente la proposition à l'utilisateur.};
\node[etape, below=0.3cm of e7] (e8) {8. Accepter / modifier : réponse mémorisée.\quad Remarques : retour à l'étape 5.};
\foreach \a/\b in {e1/e2,e2/e3,e3/e4,e4/e5,e5/e6,e6/e7,e7/e8} \draw[arr] (\a) -- (\b);
\draw[arr] (e8.east) -- ++(0.6,0) |- node[right, pos=0.25, align=left]{\footnotesize nouveau\\ \footnotesize brouillon} (e5.east);
\end{tikzpicture}
\caption{Traitement d'un email, de la réception à la réponse validée.}
\label{fig:sequence}
\end{figure}

Si l'étudiant relance ensuite dans le même fil, le \code{ThreadStep} reconnaît la
conversation, et le \code{StudentAgent} reçoit l'échange précédent dans son prompt.

Pour un email ambigu, par exemple une personne inconnue qui demande un formulaire
d'inscription pour un workshop du projet, les votes divergent : mots-clés administratifs
d'un côté, domaine universitaire de l'autre. Le routeur consulte alors l'arbitre LLM. Si
l'email contient aussi une demande relevant du projet de recherche, le
\code{ResearchProjectAgent} peut contribuer à la réponse rédigée par l'agent principal
(section~\ref{sec:collab}).

% ------------------------------------------------------------------
\section{Intégration du LLM local}

L'application communique avec un serveur local compatible avec l'API OpenAI
(\code{POST /v1/chat/completions}). Elle est prévue pour LM Studio sur le port 1234.
L'adresse et le nom du modèle sont lus dans \code{config.json}, ce qui permet de changer
de modèle sans recompiler :

\begin{center}
\small
\code{"url": "http://127.0.0.1:1234/v1/chat/completions",}\\
\code{"model": "google/gemma-4-e4b"}
\end{center}

Le reste du code ne dépend que de l'interface \code{LLM} et de sa méthode
\code{ask(systemPrompt, message)}. \code{LmStudioClient} construit la requête JSON
(un message système et un message utilisateur) et lit la réponse. En cas de problème
(serveur arrêté, modèle absent), il lève une \code{LlmException} explicite : une erreur ne
peut jamais être prise pour une réponse. Trois types d'appels sont faits au modèle :
la rédaction (agent spécialisé), l'arbitrage (classifieur) et la relecture (critique).
Si un modèle d'embeddings est configuré, une autre route est aussi utilisée,
\code{POST /v1/embeddings}, pour la sélection du contexte (section~\ref{sec:contexte}).
Les réponses JSON sont lues de façon tolérante, car les petits modèles ajoutent souvent du
texte autour.

Comme le code ne dépend que du protocole, nous avons pu utiliser aussi le serveur de
\emph{llama.cpp} (\code{llama-server}) sans modifier une ligne de Java
(section~\ref{sec:tests}). LM Studio ne fonctionne pas sur les Mac à processeur Intel.

% ------------------------------------------------------------------
\section{Sélection du contexte et journal des décisions}
\label{sec:extensions}

La première version du rapport proposait trois pistes : ne récupérer que les faits utiles
pour l'email (RAG), journaliser les décisions et comparer, sur un jeu de test, une
architecture mono-agent et l'architecture multi-agent. Ces trois pistes ont été
réalisées. Le but est de mesurer la question posée par le sujet au lieu de seulement
l'argumenter. Cette section décrit les deux premières ; la comparaison est présentée
dans la section~\ref{sec:tests}.

\subsection{Sélection du contexte envoyé au LLM}
\label{sec:contexte}

\paragraph{Problème.} L'agent choisi envoyait au modèle \emph{tous} ses faits de contexte
et \emph{tout} l'historique du fil. Les prompts grossissent avec le contexte, ce qui gêne
un petit modèle.

\paragraph{Solution.} Une étape de sélection a été ajoutée : chaque agent ne transmet que
les croyances les plus pertinentes pour l'email traité. Trois stratégies sont
interchangeables grâce au pattern \emph{Strategy} (interface \code{ContextSelector}).
Elles sont créées par une \emph{Factory} (\code{ContextSelectorFactory}).

\begin{table}[H]
\centering
\small
\begin{tabularx}{\linewidth}{l L}
\toprule
\textbf{Stratégie} & \textbf{Fonctionnement} \\
\midrule
\code{all} & Comportement d'origine : tout le contexte est envoyé. \\
\addlinespace
\code{lexical} (par défaut) & Score BM25 entre l'email et chaque fait ($k_1 = 1{,}2$ et
$b = 0{,}75$). Les textes sont d'abord normalisés : accents retirés, mots vides français
et anglais supprimés, mots réduits à une racine de 4 lettres pour rapprocher les formes
françaises et anglaises. Seuls les $k$ meilleurs faits de score non nul sont gardés,
dans leur ordre d'origine. \\
\addlinespace
\code{embedding} & Similarité cosinus entre l'email et chaque fait, sur des vecteurs
calculés par le serveur local (route \code{/v1/embeddings} de LM Studio, modèle
\emph{nomic-embed-text-v1.5}). Un décorateur (\code{CachingEmbedder}) garde les vecteurs
en mémoire pour ne pas recalculer les faits à chaque email. Si le service d'embeddings ne
répond pas, l'agent revient automatiquement à la stratégie lexicale, et ce repli est
affiché. \\
\bottomrule
\end{tabularx}
\caption{Stratégies de sélection du contexte.}
\end{table}

\paragraph{Configuration.}
\begin{itemize}
  \item Dans \code{config.json}, un bloc \code{context} règle la stratégie,
        \code{maxBeliefs} (8 par défaut) et \code{maxThreadMessages} (6 par défaut, ce
        sont les derniers messages du fil qui sont gardés).
  \item Un bloc \code{embedding} donne le nom du modèle d'embeddings. S'il est vide, la
        stratégie \code{embedding} est remplacée par \code{lexical}, avec un avertissement.
  \item Chaque agent peut redéfinir ces valeurs dans son \code{agent.json}.
\end{itemize}

\paragraph{Visibilité.} L'étape « contexte » de l'interface affiche maintenant ce qui a
été retenu. Exemple : « 3/10 context belief(s) selected (lexical), 2/2 message(s) from
this thread's memory, 450 chars ».

\paragraph{CriticAgent.} Le critique reçoit le même contexte sélectionné que l'agent
auteur, ce qui garde sa relecture cohérente avec ce que l'auteur a vu.

\paragraph{Défaut corrigé.} Pour un email sans fil connu, le contexte de base d'un agent
était ajouté deux fois au prompt (méthode \code{BeliefBase.get}). Le défaut est corrigé et
un test vérifie maintenant que chaque fait n'apparaît qu'une fois.

\subsection{Journal des décisions et explicabilité}
\label{sec:journal}

\paragraph{Principe.} Un observateur, \code{DecisionLogger}, est branché sur l'interface
\code{WorkflowListener} existante (pattern \emph{Observer}). Le routeur n'a donc pas besoin
de le connaître. Pour chaque email traité, il écrit une ligne JSON dans
\code{logs/decisions.jsonl}.

\paragraph{Contenu d'une ligne.}
\begin{itemize}
  \item \textbf{Routage} : agent choisi, rôle, confiance, score de chaque agent, signaux
        qui ont voté, arbitre consulté ou non, repli sur le \code{GenericAgent} ou non,
        explication lisible et durée.
  \item \textbf{Contexte} : stratégie, nombre de faits retenus sur le nombre disponible,
        nombre de messages du fil envoyés, taille en caractères.
  \item \textbf{Chaque brouillon} : numéro d'essai, temps de génération, taille, note et
        remarques du critique, retour de l'utilisateur et raison d'un éventuel refus.
  \item \textbf{Issue finale} (acceptée, modifiée, refusée, abandonnée) et nombre de
        brouillons.
  \item \textbf{En cas d'erreur} : l'étape qui a échoué et le message.
\end{itemize}

\paragraph{Confidentialité.} Par défaut, le texte des emails, des brouillons et des faits
n'est pas enregistré, seulement des métadonnées. L'option \code{includeText} de
\code{config.json} permet de l'activer. Le journal peut aussi être désactivé.

\paragraph{Rapport.} La commande \code{Evaluation log} relit le journal et produit un
rapport Markdown :
\begin{itemize}
  \item statistiques par agent ;
  \item taux d'arbitrage et de repli ;
  \item nombre moyen de brouillons ;
  \item calibration du critique, c'est-à-dire sa note moyenne selon la décision finale de
        l'utilisateur, pour vérifier que le critique et l'utilisateur sont d'accord ;
  \item raisons de refus les plus fréquentes.
\end{itemize}

\subsection{Nouvelles classes}

\begingroup\raggedright
\begin{itemize}
  \item \textbf{Sélection du contexte} : \code{ContextSelector}, \code{AllContextSelector},
        \code{LexicalContextSelector}, \code{EmbeddingContextSelector},
        \code{ContextSelectorFactory}, \code{ContextSelection}, \code{Terms}
        (normalisation du texte), \code{Embedder}, \code{LmStudioEmbedder},
        \code{CachingEmbedder}.
  \item \textbf{Journal} : \code{DecisionLogger}, \code{DecisionLogReport},
        \code{RoutingOutcome} (décision de routage avec scores, signaux et explication).
  \item \textbf{Évaluation} : \code{Evaluation} (point d'entrée), \code{RoutingEvaluation},
        \code{ReplyEvaluation}, \code{LabelledEmail}, \code{MeteredLLM}.
  \item \textbf{Configuration} : \code{AppConfig}, qui lit \code{config.json} et assemble
        le système.
  \item \textbf{Données} : \code{eval/dataset.json}, \code{eval/single\_agent\_prompt.txt}
        (prompt de l'agent unique), \code{eval/RESULTATS.md} (résumé des résultats).
\end{itemize}
\par\endgroup

% ------------------------------------------------------------------
\section{Collaboration entre agents sur un email ambigu}
\label{sec:collab}

\paragraph{Problème.} Jusqu'ici, un email qui concerne plusieurs agents (par exemple une
demande de formulaire d'inscription à un workshop, qui relève de l'administration
\emph{et} du projet de recherche) était confié à un seul agent, qui ignorait l'autre
partie de la demande. Le routeur peut maintenant faire répondre plusieurs agents, chacun
sur sa partie, puis fusionner leurs contributions.

\paragraph{Principe.}
\begin{enumerate}
  \item \textbf{Détection (règles).} Un agent est \emph{concerné} s'il a voté ou si l'un de
        ses mots-clés apparaît dans l'email. La collaboration n'est envisagée que si au
        moins deux agents sont concernés et que les signaux ne suffisent pas (arbitre
        consulté) : un email clair ne coûte aucun appel supplémentaire.
  \item \textbf{Répartition (LLM).} L'arbitre choisit l'agent \emph{principal} et, seulement
        si l'email contient une demande distincte, des \emph{contributeurs} avec la partie
        qui revient à chacun. Le \code{RoutingAgent} filtre cette proposition : agents
        concernés uniquement, différents de l'agent principal et du \code{GenericAgent},
        partie décrite, sans doublon, dans la limite \code{maxAgents}.
  \item \textbf{Contributions.} Nouvelle étape « Collaboration » du workflow : chaque
        contributeur rédige une courte note factuelle sur sa partie, avec \emph{ses} règles
        et \emph{son} contexte (\code{Agent.contribute}, objet \code{Contribution}).
  \item \textbf{Fusion.} L'agent principal rédige une réponse unique, avec son ton, en
        intégrant les notes. Le \code{CriticAgent} les reçoit aussi, pour ne pas compter
        comme inventée une information venue d'un contributeur.
  \item \textbf{Validation.} Inchangée. Après un refus, seul l'agent principal rédige de
        nouveau et les contributions sont réutilisées. Seul l'agent principal mémorise le
        fil.
\end{enumerate}

Les agents spécialisés ne communiquent toujours pas entre eux : les notes passent par le
\code{RoutingAgent} (médiateur) et l'\code{Intention} de l'agent principal. Un contributeur
en panne n'arrête pas le traitement : l'étape est marquée en échec et l'agent principal
répond seul. Si le modèle ne propose aucun contributeur, le comportement est celui d'avant.

\paragraph{Configuration.} Dans \code{config.json} :
\code{"collaboration": \{"enabled": true, "maxAgents": 2\}}. \code{maxAgents} compte
l'agent principal (2 = un contributeur au plus) ; \code{enabled: false} désactive le
mécanisme.

\paragraph{Visibilité.}
\begin{itemize}
  \item Interface : onzième étape « Collaboration » dans le suivi du workflow ; décision de
        routage avec agent principal et contributeurs (graphe des interactions et pastilles) ;
        bandeau « Brouillon rédigé par X avec la contribution de Y (partie) » et bouton
        \emph{Contributions} pour lire chaque note.
  \item Journal des décisions : contributeurs, partie, taille, durée et erreur éventuelle ;
        le rapport indique le taux de collaboration.
\end{itemize}

\paragraph{Évaluation.} Deux configurations de comparaison ont été ajoutées au banc :
\code{no-collaboration} pour le routage (taux de collaboration, part des collaborations
pertinentes, couverture des emails ambigus) et \code{multi-agent-no-collab} pour la
qualité des réponses, à comparer à \code{multi-agent} sur les emails ambigus. Pour ces
emails, le juge connaît aussi le contexte des autres agents acceptables. L'email de
démonstration~5 illustre le cas : \code{AdministrationAgent} rédige avec la contribution
de \code{ResearchProjectAgent}.

\paragraph{Tests et limites.} Une nouvelle section de tests (21 vérifications, avec LLM
simulé) couvre la répartition, le filtrage, la fusion, la critique, le nouveau brouillon,
la mémoire, la panne d'un contributeur, la désactivation et le journal ; la suite passe
de 78 à 99 vérifications. Ces tests valident la logique, pas la qualité des réponses : la
collaboration n'a pas encore été mesurée avec un vrai modèle, et un petit modèle peut ne
pas remplir la liste des contributeurs (l'agent principal répond alors seul).

\paragraph{Nouvelles classes et modifications.}
\begingroup\raggedright
\begin{itemize}
  \item \textbf{Nouvelle classe} : \code{Contribution}.
  \item \textbf{Routage} : \code{RoutingDecision} (contributeurs), \code{Arbiter} et
        \code{LlmClassifierStep} (désignation des contributeurs), \code{RoutingStep} et
        \code{KeywordStep} (agents concernés), \code{RoutingAgent} (filtrage, étape de
        collaboration).
  \item \textbf{Agents} : \code{Agent}, \code{BDIAgent} (\code{contribute}, fusion dans le
        brouillon), \code{Intention}, \code{Proposal}, \code{CriticAgent}.
  \item \textbf{Workflow et interface} : \code{WorkflowStage}, \code{WorkflowListener},
        \code{WorkflowPanel}, \code{RoutingPanel}, \code{ProposalPanel},
        \code{AgentNetworkView}, \code{ProcessPage}, \code{DashboardPage},
        \code{ConsoleValidator}.
  \item \textbf{Journal, évaluation et configuration} : \code{DecisionLogger},
        \code{DecisionLogReport}, \code{RoutingOutcome}, \code{RoutingEvaluation},
        \code{ReplyEvaluation}, \code{AppConfig}, \code{config.json}.
\end{itemize}
\par\endgroup

% ------------------------------------------------------------------
\section{Interface graphique}

L'interface est développée en Swing et comporte quatre pages. Le traitement d'un email
tourne dans un thread séparé, pour que la fenêtre ne se fige pas pendant les appels au
modèle. Quand le routeur attend la décision de l'utilisateur, \code{SwingValidator}
bloque ce thread jusqu'au clic sur un bouton.

\begin{itemize}
  \item \textbf{Dashboard} : agents chargés, état du LLM, statistiques et activité de la session.
  \item \textbf{Process email} : saisie ou chargement d'un email de démonstration, suivi des
        onze étapes du traitement (\emph{Pending}, \emph{Processing}, \emph{Completed},
        \emph{Failed}, \emph{Skipped}), décision de routage expliquée (agent principal et
        contributeurs) et proposition à valider (bouton \emph{Contributions}).
  \item \textbf{Agents} : consultation et modification de la description, des mots-clés, du
        \emph{system prompt} et du contexte de chaque agent, consultation de la mémoire et
        création d'un nouvel agent. Les modifications sont enregistrées dans les fichiers
        et prises en compte immédiatement.
  \item \textbf{Local LLM} : adresse du serveur, modèle configuré, état de la connexion,
        prompt de test et activation du CriticAgent.
\end{itemize}

\capture{dashboard.png}{Tableau de bord : agents disponibles, état du LLM local et activité.}{fig:dashboard}
\capture{routing.png}{Décision de routage : agent choisi, confiance, signaux et agents impliqués.}{fig:routing}
\capture{workflow.png}{Suivi en direct des étapes du traitement d'un email.}{fig:workflow}
\capture{proposal.png}{Proposition générée, avis du CriticAgent et boutons de validation.}{fig:proposal}
\capture{agents.png}{Gestion des agents : modification du \emph{system prompt} et du contexte.}{fig:agents}

% ------------------------------------------------------------------
\section{Tests et expérimentations}
\label{sec:tests}

Nous avons préparé six emails de démonstration (\code{demo/emails.json}) qui couvrent les
cas demandés par le sujet. Trois types de vérifications ont été faites, que nous
distinguons clairement. Un banc d'expériences plus large, sur 30 emails étiquetés, est
présenté à la fin de cette section (section~\ref{sec:banc}).

\begin{table}[H]
\centering
\small
\begin{tabularx}{\linewidth}{l L l}
\toprule
\textbf{Email} & \textbf{Signaux attendus} & \textbf{Agent} \\
\midrule
1. Étudiant (examen) & contact + mot-clé & StudentAgent \\
2. Collègue (réunion de jury) & contact + mots-clés & TeacherAgent \\
3. Scolarité (documents, échéance) & contact + mots-clés & AdministrationAgent \\
4. Coordinatrice SuperMASProject & contact + mots-clés & ResearchProjectAgent \\
5. Email ambigu (workshop) & conflit Administration / Teacher, arbitre LLM & selon l'arbitre \\
6. Relance de l'email 1 & fil déjà traité & StudentAgent \\
\bottomrule
\end{tabularx}
\caption{Emails de démonstration.}
\end{table}

\paragraph{1. Tests automatisés (sans vrai modèle).}
La commande \code{make test} lance 99 vérifications avec un LLM simulé et un utilisateur
simulé : routage des cinq premiers emails, arbitrage, repli sur le \code{GenericAgent},
acceptation, modification, nouveau brouillon, refus, abandon, mémoire du fil,
CriticAgent, ordre des étapes, panne du LLM, création et modification d'un agent.
De nouveaux scénarios couvrent la sélection du contexte, le journal des décisions, le
rapport du journal, le banc d'expériences, la correction du contexte dupliqué et la
collaboration entre agents : la suite est passée de 41 à 99 vérifications. Les 99
vérifications passent. Ces tests valident la logique multi-agent, pas la qualité
des réponses.

\paragraph{2. Tests avec un serveur simulé.}
L'interface a été parcourue de bout en bout avec un petit serveur local qui implémente
les mêmes routes HTTP que LM Studio mais renvoie des textes fixes. Les six emails ont été
routés vers les agents attendus, et l'email 5 a bien déclenché l'arbitre. La
modification, le refus, la régénération avec remarques et la mémoire du fil ont fonctionné.
Le cas « serveur arrêté » affiche une erreur claire, et rien n'est mémorisé.
Ces essais valident le chemin HTTP et l'interface, \emph{pas} un vrai modèle.

\paragraph{3. Essai avec un vrai modèle local.}
LM Studio n'étant pas disponible sur la machine de test (Mac Intel), nous avons utilisé
\code{llama-server} (llama.cpp) avec le modèle \emph{Qwen2.5-3B-Instruct} (quantification
Q4\_K\_M), uniquement sur processeur. Pour l'email 1 :
\begin{itemize}
  \item le routage a choisi le \code{StudentAgent} avec une confiance de 1,00 (contact + mot-clé), sans appel au LLM ;
  \item le brouillon a été généré en 27 s environ et la relecture du critique en 42 s environ ;
  \item le CriticAgent a donné 4/10 avec des remarques pertinentes : la question sur les
        chapitres n'était pas traitée et les échéances n'étaient pas rappelées.
\end{itemize}

\paragraph{Observations.}
Le routage par signaux est rapide et explicable, et il évite un appel au modèle dans les
cas clairs. Le critique repère des oublis réels du petit modèle, ce qui justifie sa
présence. En revanche, chaque appel coûte plusieurs dizaines de secondes sur
processeur : le critique double le temps de traitement. Nous avons aussi constaté qu'un
modèle de 3B paramètres peut affirmer des choses fausses sur lui-même, ce qui confirme
qu'il faut toujours une validation humaine.

\subsection{Banc d'expériences}
\label{sec:banc}

\paragraph{Jeu de données.} Le fichier \code{eval/dataset.json} contient 30 emails
étiquetés. Chaque email indique l'agent attendu et, si besoin, d'autres agents
acceptables. Les emails sont répartis en cinq catégories : cas clairs, expéditeurs
inconnus, emails en français, contexte croisé (un contact connu qui écrit sur un sujet
d'un autre contexte) et emails ambigus. Quatorze emails sont marqués pour l'évaluation de
la qualité des réponses.

\paragraph{Commande.} \code{java -cp "bin;lib/*" Evaluation routing|replies|log}, ou les
cibles \code{make eval-routing}, \code{make eval-replies} et \code{make eval-log}. Les
rapports Markdown, CSV et JSON sont écrits dans \code{eval/results/}. Un décorateur,
\code{MeteredLLM}, compte les appels au modèle, la taille des prompts et le temps passé.

\paragraph{Expérience 1 : routage.} Cinq configurations du routeur sont comparées.

\begin{table}[H]
\centering
\small
\begin{tabularx}{\linewidth}{l L}
\toprule
\textbf{Configuration} & \textbf{Description} \\
\midrule
full & signaux + arbitre LLM (système complet) \\
rules-only & signaux seuls, sans arbitre \\
llm-only & arbitre LLM seul, sans signaux \\
no-keywords & sans le signal mots-clés \\
content-only & sans les signaux contact et fil (contenu seul) \\
\bottomrule
\end{tabularx}
\caption{Configurations du routeur comparées.}
\end{table}

Le rapport donne l'exactitude stricte (agent attendu) et l'exactitude tolérante (agent
acceptable), les taux d'arbitrage et de repli, la matrice de confusion, la précision, le
rappel et le F1 par agent, l'exactitude par catégorie et la liste des erreurs.

\paragraph{Expérience 2 : qualité des réponses.} Quatre configurations rédigent une
réponse pour chaque email.

\begin{table}[H]
\centering
\small
\begin{tabularx}{\linewidth}{l L}
\toprule
\textbf{Configuration} & \textbf{Description} \\
\midrule
multi-agent & routage + agent spécialisé, contexte sélectionné (BM25) \\
multi-agent, tout le contexte & routage + agent spécialisé, sans sélection \\
oracle & l'agent attendu est imposé, pour isoler la rédaction des erreurs de routage \\
agent unique & baseline : un prompt générique et le contexte de \emph{tous} les agents, sans routage \\
\bottomrule
\end{tabularx}
\caption{Configurations comparées pour la qualité des réponses.}
\end{table}

Chaque brouillon est noté par le \code{CriticAgent}, qui joue le rôle de juge, avec les
règles et le contexte de l'agent attendu. Le rapport donne la note moyenne, le taux de
brouillons approuvés ($\geq$ 7/10), la taille du contexte et du prompt, le temps par
brouillon et la note par email. L'option \code{--judge-model} permet de prendre un modèle
juge différent du rédacteur.

\subsection{Résultats avec LM Studio}
\label{sec:resultats}

\paragraph{Conditions.} Serveur LM Studio local ; le modèle \emph{google/gemma-4-e4b}
rédige, arbitre et juge. Les embeddings sont calculés par \emph{nomic-embed-text-v1.5}.
Un appel au modèle prend environ 25\,s. C'est le premier essai du système avec LM Studio
lui-même ; les essais précédents utilisaient llama.cpp.

\begin{table}[H]
\centering
\small
\begin{tabularx}{\linewidth}{L c c c}
\toprule
\textbf{Configuration} & \textbf{Exactitude stricte} & \textbf{Exactitude tolérante} & \textbf{Appels LLM} \\
\midrule
full (signaux + arbitre) & 87\,\% & 93\,\% & 15 \\
rules-only & 63\,\% & 63\,\% & 0 \\
llm-only & 87\,\% & 97\,\% & 30 \\
no-keywords & 80\,\% & 90\,\% & -- \\
content-only & 90\,\% & 97\,\% & -- \\
\bottomrule
\end{tabularx}
\caption{Routage des 30 emails étiquetés.}
\label{tab:routage}
\end{table}

\paragraph{Routage.}
\begin{itemize}
  \item L'arbitre LLM apporte 24 points par rapport aux signaux seuls.
  \item La configuration \emph{full} atteint la même exactitude stricte que le LLM seul
        avec deux fois moins d'appels au modèle : les signaux règlent les cas clairs sans
        LLM.
  \item Deux erreurs de \emph{full} sont instructives. Un collègue connu du projet de
        recherche qui écrit au sujet d'un cours est envoyé au
        \code{ResearchProjectAgent} : le vote du contact (1,0) suffit seul et l'arbitre
        n'est pas consulté. Un « rapport d'activité annuel » en français part vers le
        \code{TeacherAgent} dans toutes les configurations, faute de mots-clés
        administratifs en français.
  \item La configuration \emph{content-only}, qui fait mieux ici, montre que le signal
        « contact » peut nuire quand le contenu contredit le rôle habituel de
        l'expéditeur.
\end{itemize}

\begin{table}[H]
\centering
\small
\begin{tabularx}{\linewidth}{L c c c c c}
\toprule
\textbf{Configuration} & \textbf{Note /10} & \textbf{Approuvés} & \textbf{Contexte (car.)} & \textbf{Prompt (car.)} & \textbf{Temps} \\
\midrule
multi-agent (sélection BM25) & 7,83 & 83\,\% & 330 & 1\,523 & 49\,s \\
multi-agent, tout le contexte & 8,00 & 80\,\% & 731 & 2\,058 & 50\,s \\
oracle (agent attendu imposé) & 8,83 & 100\,\% & 243 & 1\,361 & 43\,s \\
agent unique & 8,25 & 100\,\% & 3\,733 & 6\,102 & 62\,s \\
\bottomrule
\end{tabularx}
\caption{Qualité des réponses sur 6 emails (temps par brouillon).}
\label{tab:reponses}
\end{table}

\paragraph{Qualité des réponses.} Ces résultats donnent un premier élément de réponse à
la question du sujet.
\begin{itemize}
  \item Quand le bon agent répond, l'organisation multi-agent obtient de meilleures notes
        que l'agent unique (8,83 contre 8,25). Son prompt est 4,5 fois plus court et ses
        réponses environ 30\,\% plus rapides.
  \item L'écart du multi-agent complet vient d'un seul email mal routé, l'erreur de
        contexte croisé ci-dessus, noté 3/10. Sans cet email, la moyenne est de 8,8. La
        qualité du système dépend donc d'abord de la qualité du routage.
  \item La sélection du contexte divise la taille du contexte par deux sans baisse de
        qualité mesurable (7,83 contre 8,00, écart non significatif sur 6 emails).
\end{itemize}

\paragraph{Limites de ces résultats.}
\begin{itemize}
  \item L'échantillon est petit : 6 emails pour les réponses.
  \item Le juge est le même petit modèle que le rédacteur, et il s'est montré indulgent :
        des brouillons qui contenaient des champs à compléter (\code{[Your Name]}) ont
        reçu 9 ou 10.
  \item Une réponse en anglais à un email en français n'a pas été pénalisée.
  \item Dans 3 cas sur 24, la réponse du juge n'a pas pu être lue. Ces cas sont exclus des
        moyennes, ce qui favorise légèrement l'agent unique.
\end{itemize}

\paragraph{Embeddings contre BM25 (essai qualitatif).} La sélection par embeddings
fonctionne (0,1 à 0,8\,s par email avec le cache, environ 10\,s au premier appel). Elle est
cependant moins pertinente que BM25 sur les emails en français. Les préfixes de requête
attendus par le modèle nomic ne sont pas encore utilisés.

\paragraph{Améliorations identifiées (non réalisées).}
\begin{itemize}
  \item Baisser le poids du signal « contact » (environ 0,6) pour qu'un contenu
        contradictoire déclenche l'arbitre.
  \item Ajouter des mots-clés administratifs en français (rapport, activité\ldots).
  \item Utiliser les préfixes attendus par le modèle d'embeddings.
  \item Imposer dans les prompts de répondre dans la langue de l'email, sans champ à
        compléter.
  \item Refaire l'évaluation sur les 14 emails avec un juge plus grand.
\end{itemize}

% ------------------------------------------------------------------
\section{Limites et perspectives}

\paragraph{Limites.}
\begin{itemize}
  \item Avec LM Studio, l'évaluation porte sur peu d'emails (30 pour le routage, 6 pour la
        qualité des réponses), et le juge est le même petit modèle que le rédacteur.
  \item Les mots-clés sont lexicaux : un synonyme ou une forme non listée ne vote pas. Il
        manque en particulier des mots-clés administratifs en français.
  \item Le vote du contact suffit seul : quand le contenu contredit le rôle habituel de
        l'expéditeur, l'arbitre n'est pas consulté.
  \item Les emails sont saisis à la main ; il n'y a pas de connexion à une vraie boîte mail.
  \item La mémoire n'envoie que les derniers messages du fil (6 par défaut), mais il n'y a
        pas de résumé pour les longs échanges : les messages plus anciens sont ignorés.
  \item Sur processeur, la latence est élevée, surtout avec le critique activé.
\end{itemize}

\paragraph{Perspectives.}
Trois pistes de la première version ont été réalisées : la sélection des faits utiles
par BM25 ou par \emph{embeddings} (RAG), le journal des décisions et la comparaison avec
un agent unique (sections~\ref{sec:extensions} et~\ref{sec:banc}). Il reste à utiliser
des \emph{embeddings} pour un routage sémantique, résumer les longs fils, ajuster les
mots-clés à partir des corrections de l'utilisateur et du journal, et refaire
l'évaluation sur plus d'emails avec un juge plus grand. Les améliorations identifiées
par les expériences sont listées en section~\ref{sec:resultats}. Une importation des
emails en lecture seule (IMAP) serait possible, à condition de conserver la validation
humaine avant tout envoi.

% ------------------------------------------------------------------
\section{Conclusion}

Ce projet nous a permis de construire un vrai système multi-agent autour d'un LLM local.
Les agents ont des rôles, des connaissances et des règles distincts. Un coordinateur
explicable décide qui répond et arbitre les cas ambigus. Le cycle BDI organise la
rédaction et la validation, et un agent critique contrôle la qualité. L'organisation
donne au modèle des prompts courts et ciblés, isole les contextes et reste facile à
étendre : un agent s'ajoute sans modifier le code, et le moteur LLM peut être remplacé
par un simple changement de configuration. La principale limite reste la qualité et la
vitesse des petits modèles locaux, que l'organisation compense en partie sans les supprimer.

Les premières mesures avec LM Studio (section~\ref{sec:resultats}) vont dans ce sens :
quand le bon agent répond, les réponses sont mieux notées qu'avec un agent unique, avec un
prompt 4,5 fois plus court. Elles montrent aussi que la qualité dépend d'abord du
routage. Ces résultats restent à confirmer sur plus d'emails et avec un juge plus grand.

% ------------------------------------------------------------------
\section{Contributions des membres du groupe}

\begin{table}[H]
\centering
\small
\begin{tabularx}{\linewidth}{l l L}
\toprule
\textbf{Étudiant} & \textbf{N° étudiant} & \textbf{Contributions} \\
\midrule
Mounia AREZZOUG & 22509499 &
Participation à la conception de l'architecture initiale (routage, agents BDI) ;
ajout de l'AdministrationAgent, du ResearchProjectAgent et du CriticAgent ;
développement de l'interface graphique Swing (avec assistance IA, voir la déclaration) ;
emails de démonstration et tests ; mise en place et essai avec un modèle local
(llama.cpp, Qwen2.5-3B) ; mise à jour du README et intégration sur GitHub ;
rédaction et mise en forme du rapport final. \\
\addlinespace
Yunus Emre DOGAN & 22508757 &
Version de base du système (agents BDI, croyances et mémoire des fils,
fabrique d'agents, client LLM, contacts) ; premiers agents (StudentAgent,
TeacherAgent, GenericAgent) avec leurs prompts et contextes ; routage par vote
de signaux avec seuils, arbitre LLM et repli GenericAgent ; boucle de
validation humaine (intention BDI, accepter / modifier / refuser,
\texttt{maxRetries}, mémorisation après validation) ; configuration du LLM et
gestion des pannes ; Makefile, README initial, dépôt GitHub. \\
\addlinespace
Mathis LEVESQUE DU ROSTU & 22510354 &
Sélection du contexte (BM25, embeddings locaux avec cache et repli) ; journal des
décisions et rapport d'explicabilité ; jeu de 30 emails étiquetés et banc d'expériences
(routage, multi-agent contre agent unique) ; exécution des expériences avec LM Studio et
analyse des résultats ; correction du contexte dupliqué et tests associés. \\
\addlinespace
Nathan THIERY & 22512978 &
Collaboration entre agents sur un email ambigu (section~\ref{sec:collab}) : conception et
implémentation (agents concernés, répartition par l'arbitre LLM, filtrage des
contributeurs, notes des contributeurs, fusion par l'agent principal, relecture du critique
avec les notes) ; adaptation de l'interface (étape « Collaboration », agent principal et
contributeurs dans le graphe, bouton \emph{Contributions}), du journal des décisions et du
banc d'expériences (configurations \code{no-collaboration} et \code{multi-agent-no-collab}) ;
paramétrage dans \code{config.json}.\\
\bottomrule
\end{tabularx}
\caption{Répartition des contributions.}
\end{table}

\section*{Déclaration d'utilisation de l'intelligence artificielle}
\addcontentsline{toc}{section}{Déclaration d'utilisation de l'intelligence artificielle}

Des outils d'intelligence artificielle ont été utilisés comme assistance pendant ce projet.
Ils ont principalement servi à améliorer la présentation visuelle de l'interface graphique,
afin d'obtenir un rendu plus soigné et agréable à utiliser. Le code de l'interface a été
développé puis amélioré progressivement, et non entièrement généré. Une assistance a
également été utilisée ponctuellement lors de certaines étapes de développement, notamment
pour l'ajout d'agents, de l'agent critique et des tests.

L'architecture de base, le routage par signaux, le modèle BDI et la mémoire des conversations
sont issus du travail du groupe. Les éléments ayant bénéficié d'une assistance par IA ont
été étudiés et compris. Nous sommes en mesure d'expliquer l'architecture du projet, le
fonctionnement des agents et les choix techniques réalisés.

Pour la contribution de Mathis Levesque du Rostu (sélection du contexte, journal des
décisions, banc d'expériences), un assistant de programmation IA (Claude Code) a servi à
proposer les extensions, à écrire une grande partie du code et des tests, et à lancer les
expériences. Les choix, les résultats et leurs limites ont été relus et analysés.

La partie tests et mise à jour du README de l'implémentation de la fonctionnalité de collaboration entre agents sur un email ambigu (section~\ref{sec:collab}) a également
été implémentée avec l'aide de Claude Code.

% ------------------------------------------------------------------
\appendix
\section{Réponses aux questions de conception du sujet}

\begin{table}[H]
\centering
\small
\begin{tabularx}{\linewidth}{>{\raggedright\arraybackslash}p{4.3cm} L}
\toprule
\textbf{Question} & \textbf{Réponse} \\
\midrule
Qu'est-ce qu'un agent ? & Une entité autonome avec un rôle, un prompt, ses croyances, un désir et des intentions (\code{BDIAgent}), plus le coordinateur et le critique. \\
Ses responsabilités ? & Routeur : choisir et coordonner. Agent spécialisé : rédiger et mémoriser. Critique : évaluer. \\
Quelles informations possède-t-il ? & Ses fichiers de contexte et la mémoire de ses propres fils. \\
Comment communiquent-ils ? & Par objets-messages (Email, Intention, Proposal, Feedback, Critique), via le RoutingAgent. \\
Comment choisir l'agent ? & Votes de quatre signaux, seuils 75\,\% et 0,7, sinon arbitre LLM. \\
Plusieurs agents conviennent ? & L'arbitre tranche entre les agents en conflit ; si l'email contient des demandes distinctes, il désigne un agent principal et des contributeurs dont les notes sont fusionnées (section~\ref{sec:collab}). \\
Aucun ne convient ? & L'arbitre choisit parmi tous ; si confiance < 0,6, GenericAgent. \\
Représentation des contextes ? & Fichiers texte \code{date | fait} par agent, modifiables depuis l'interface. \\
Historique des conversations ? & Croyances EMAIL / REPLY par fil, enregistrées après validation seulement. \\
Éviter trop de contexte ? & Seuls le contexte de l'agent choisi et le fil courant sont utilisés ; parmi eux, seuls les faits les plus pertinents (BM25 ou embeddings, 8 au plus) et les 6 derniers messages du fil sont envoyés. \\
Ajouter un agent ? & Créer un dossier (ou utiliser l'interface), sans modifier le code. \\
Changer de LLM ? & Modifier \code{config.json} ; un autre moteur = une classe qui implémente \code{LLM}. \\
Données persistantes ? & \code{agent.json}, prompts, contextes, contacts, mémoire validée, \code{config.json}. \\
Avantages et limites ? & Prompts ciblés, routage explicable, isolation, contrôle humain ; appels LLM supplémentaires et dépendance aux mots-clés. \\
\bottomrule
\end{tabularx}
\end{table}

\end{document}